\section{Преглед на съществуващите решения}

За да бъде правилно оценена една софтуерна разработка, тя следва да бъде позиционирана спрямо вече съществуващите решения в областта. Това е особено важно при системите за електронна търговия, тъй като пазарът предлага както готови платформи с богата функционалност, така и множество технологични стекове за изграждане на custom приложения. Настоящият проект не възниква във вакуум; той е реализиран в контекста на вече наложени модели за организация на каталог, поръчка, потребителски профил и административни функции. Целта на настоящата глава е да очертае този контекст и да покаже с какво разглежданата система се доближава до утвърдените практики и в какво съзнателно се различава.

\subsection{Обща картина на съвременните решения за електронна търговия}

Съвременните системи за електронна търговия могат условно да се разделят на три основни групи. Първата група включва напълно управлявани SaaS платформи, при които доставчикът поема голяма част от инфраструктурните и административните задачи. Втората група включва self-hosted системи с готова функционалност и възможност за разширяване чрез плъгини и теми. Третата група обхваща custom решения, разработени върху общопредназначени уеб рамки, при които логиката, структурата на данните и интерфейсът се изграждат специално за конкретния проект.

Изборът между тези подходи зависи от множество фактори: степен на необходима гъвкавост, бюджет, срокове, изисквания към интеграции, организационен контекст и налични технически компетентности. Готовите платформи предлагат бърз старт и богат набор от стандартни функции, но ограничават архитектурния контрол. Custom разработката изисква повече проектантски и програмни усилия, но предоставя възможност системата да бъде моделирана точно според целите на екипа или учебната задача. В този смисъл настоящият проект попада в третата група и следва да бъде разглеждан като контролирана инженерна реализация, а не като конфигуриране на готов продукт.

\subsection{Готови платформи: Shopify, WooCommerce и Magento}

Една от най-разпознаваемите SaaS платформи за електронна търговия е Shopify \cite{shopify}. Нейното основно предимство е, че предлага цялостна инфраструктура, готов административен панел, управление на продукти, теми, плащания и интеграции с външни услуги. За бизнеси, които търсят бързо излизане на пазара и не се нуждаят от ниско ниво на контрол над архитектурата, Shopify е разумен избор. От гледна точка на дипломна разработка обаче този тип платформа е по-слабо подходяща, защото голяма част от инженерните решения са вече предварително взети от доставчика, а възможностите за задълбочен архитектурен анализ са ограничени.

WooCommerce \cite{woocommerce} представлява различен подход. То е разширение за WordPress, което превръща система за управление на съдържание в електронен магазин. Предимството му е в лесната инсталация, голямата екосистема от теми и плъгини и ниската бариера за въвеждане. Подходът е особено полезен, когато вече съществува WordPress сайт и електронният магазин трябва да бъде интегриран към него. Недостатъкът е, че архитектурата на приложението остава силно зависима от структурата на WordPress и от начина, по който плъгините взаимодействат помежду си. Това ограничава чистотата на модела и прозрачността на вътрешната бизнес логика.

Magento Open Source \cite{magento} е пример за мощна self-hosted платформа, ориентирана към по-сложни сценарии, по-големи каталози и по-високи изисквания към конфигурация и разширяване. Тя предлага богат набор от модули, административни възможности и интеграционни точки. В същото време сложността ѝ е значително по-висока от тази на WooCommerce или SaaS решенията. За малка или учебна разработка такава платформа често е прекалено тежка, тъй като значителна част от усилието се пренасочва към усвояване на рамката и нейните вътрешни механизми, вместо към ясно проектно моделиране на конкретното приложение.

Съпоставката между тези платформи показва, че готовите решения са силни, когато бизнесът има нужда от бърза operational готовност. Те обаче не са оптимални, когато основната цел е демонстрация и анализ на фундаментални принципи от софтуерното инженерство. Именно поради това в настоящия проект е предпочетена самостоятелна разработка върху обща уеб рамка.

\begin{table}[H]
    \centering
    \caption{Съпоставка между типични подходи за изграждане на електронен магазин}
    \label{tab:ecommerce-platforms}
    \begin{tabularx}{\textwidth}{L{2.7cm}YYY}
        \toprule
        Подход & Предимства & Ограничения & Подходящост за дипломна разработка \\
        \midrule
        Shopify & Бърз старт, хоствана инфраструктура, готови интеграции & Ограничен архитектурен контрол, зависимост от доставчик & Ниска, защото инженерните решения са силно абстрахирани \\
        WooCommerce & Лесно внедряване, богата екосистема, ниска бариера за вход & Силна зависимост от WordPress, натрупване на плъгини & Средна, ако фокусът е върху конфигурация, не върху архитектура \\
        Magento & Богата функционалност, разширяемост, зрял търговски модел & Висока сложност и тежка среда за внедряване & Средна към ниска при ограничен времеви ресурс \\
        Custom Spring приложение & Пълен контрол върху модел, код и архитектура & По-голяма първоначална разработка & Висока, защото позволява анализ на всички слоеве \\
        \bottomrule
    \end{tabularx}
\end{table}

Таблица \ref{tab:ecommerce-platforms} показва, че custom Spring приложение не е най-бързият път към готова търговска платформа, но е най-подходящият избор, когато целта е демонстрация на архитектурна последователност, моделиране на предметна област и проследимост на решенията от потребителския интерфейс до базата данни.

\subsection{Custom уеб приложения и ролята на Spring екосистемата}

В Java екосистемата Spring Boot е сред най-често използваните средства за бързо изграждане на самостоятелни уеб приложения \cite{springboot}. Неговото основно предимство се състои в стандартизирания модел за конфигурация, автоматична инициализация на инфраструктурата и лесна интеграция между различни модули като MVC, сигурност, достъп до данни, шаблони и външни услуги. Това го прави подходяща основа за приложения, които трябва да имат ясен архитектурен гръбнак, но да не се обременяват от прекомерна ръчна конфигурация.

Spring MVC \cite{springmvc} е особено важен в контекста на настоящата разработка, защото предлага естествен модел за маршрутизация на HTTP заявки, binding на параметри, интеграция с валидация и лесна връзка с шаблонни енджини като Thymeleaf \cite{thymeleaf}. По този начин всеки URL маршрут може ясно да бъде свързан с контролер, бизнес операция и визуализация. Това е съществено предимство при дипломна разработка, защото позволява по-ясно описание на потоците в системата.

Spring Data JPA \cite{springdatajpa} също играе важна роля. Вместо значителна част от времето да бъде изразходвана за ръчно писане на SQL слой и трансформации между резултатни множества и обекти, JPA моделът позволява домейн обектите да бъдат описани като entity класове с релации, а repository интерфейсите да делегират част от рутинната инфраструктурна логика на рамката. Това не премахва нуждата от добро моделиране, но ускорява реализацията на устойчив слой за съхранение на данни.

За малки и средни приложения изборът на Spring Boot носи и друго предимство: позволява решението да остане монолитно в архитектурен смисъл, но вътрешно добре разделено по модули. Това е особено подходящо за учебен проект, при който е по-важно вътрешната организация да бъде прозрачна и последователна, отколкото приложението да бъде преждевременно разделяно в множество услуги.

\subsection{Архитектурни подходи при уеб приложенията}

Освен между различни платформи, изборът трябва да бъде направен и между различни архитектурни парадигми на самото уеб приложение. В най-общ план могат да бъдат разграничени два доминиращи модела. Първият е сървърно рендираното приложение, при което HTML се генерира на сървъра и се изпраща към клиента като готова страница. Вторият е клиентски ориентираният модел, при който front-end приложението е самостоятелно, а сървърът предоставя основно REST или GraphQL API.

Сървърно рендираният модел има няколко съществени предимства за разглеждания проект. Той намалява общата архитектурна сложност, тъй като представящият слой и бизнес логиката остават в една и съща технологична среда. Освен това позволява бърза реализация на форми, валидации и навигационни потоци, без да е необходимо изграждането на отделно front-end приложение с независим процес на компилация и разгръщане. За дипломен проект, чиято основна цел е да покаже цялостен back-end и UI процес с ограничен технологичен риск, това е рационален избор.

Клиентски ориентираният модел, типичен за приложения, базирани на React, Angular или Vue, предоставя по-голяма динамичност на интерфейса и по-естествена интеграция с мобилни и външни клиенти. Той обаче изисква отделна front-end архитектура, по-ясно API контрактуване, повече внимание към клиентското състояние и често по-сложен модел за сигурност и автентикация. В контекста на текущата задача това би довело до разширяване на технологичния обхват за сметка на фокуса върху систематичната реализация на основните бизнес процеси.

Настоящият проект избира хибриден, но доминиращо сървърно рендиран подход. Основната функционалност се визуализира чрез Thymeleaf, докато отделна REST крайна точка е използвана за демонстрация на програмно достъпна бизнес информация. По този начин системата запазва простота, но същевременно показва възможност за комбиниране на HTML и JSON интерфейси.

\subsection{Позициониране на настоящия проект}

Разработваната система следва да бъде разбрана като custom монолитно приложение, ориентирано към учебно-приложна стойност. То не се конкурира пряко с пълнофункционални търговски платформи, а по-скоро демонстрира как чрез добре подбрани технологии могат да бъдат реализирани основните механизми на електронен магазин в прозрачна и разбираема архитектурна форма.

Проектът заема междинна позиция между два екстрема. От едната страна са платформите тип „всичко готово“, които скриват значителна част от инженерния процес. От другата страна са големите enterprise решения с микросервизна или силно модулна архитектура, чиято сложност е непропорционална на целите на дипломната задача. Настоящата система избира балансиран път: достатъчно реалистична функционалност, реална база данни, реални HTML шаблони, реална e-mail интеграция и реална структура на сорс кода, но без прекомерни инфраструктурни усложнения.

Тази позиция я прави особено подходяща за академичен анализ. В нея могат ясно да бъдат проследени домейн обектите, заявките към системата, маршрутизацията, ролята на валидацията, управлението на сесията, трансформациите между модели и мястото на конфигурацията. Именно тази проследимост е едно от основните предимства на избрания подход.

\subsection{Изводи от обзорната глава}

Прегледът на съществуващите решения показва, че в областта на електронната търговия няма универсално най-добър подход; подходящият избор зависи от целите на конкретния проект. Ако основната задача е бързо въвеждане в експлоатация, SaaS и плъгин-базираните решения имат ясни предимства. Ако обаче водеща е необходимостта от архитектурна прозрачност, контрол върху бизнес логиката и възможност за детайлно инженерно описание, custom уеб приложение върху Spring екосистемата е по-подходящ избор.

Именно в този смисъл настоящата дипломна разработка е обоснована: тя не се стреми да възпроизведе в пълнота зрелите търговски платформи, а да изгради представителен, работещ и аналитично обясним модел на онлайн магазин. Следващите глави развиват тази идея, като първо въвеждат необходимите теоретични основи, а след това последователно представят архитектурата и реализацията на конкретната система.
